

%====================
\section{Plasticity Implementation}
%====================
\subsection{Incremental Initial Value 
Problem for Elastoplasticity with Isotropic Hardening}

To motivate the time stepping components of the 
implementation algorithm, consider
smooth function $f(x): \real \rightarrow \real$.
With initial condition $x(0) = x_0$ 
and the notation $x_n \approx x(t_n)$, $f$ can be 
approximated as 
\begin{align}
 \dot{x}(t) & = f\left(x(t)\right), \\
\approx  \dst\frac{x_{n+1} - x_n}{\Delta t} & = f(x_{n+\theta}), \;\; 
\mbox{where} \;\; \theta \in [0, 1], \;\; \mbox{and}
 \label{eq:timeStepApproximation} \\
  x_{n+\theta} & = (1-\theta) x_n + \theta x_{n+1}.
\end{align}
The value $\theta$ controls the stability and 
accuracy of the time stepping algorithm.  
\be
\begin{tabular}{ll}
 $\theta = 0$ & forward (explicit) Euler, \\
 $\theta = \frac{1}{2}$ & midpoint rule,  \\
 $\theta = 1$ & backward (implicit) Euler.
\end{tabular}
\ee
With these notions at hand and use of the implicit Euler, 
we use~(\ref{eq:timeStepApproximation}) for evolution in the plastic flow rule.
\begin{align}
\hstrainpdot \approx \dst\frac{\hstrainpnp - \hstrainpn}{\Delta t} 
& = \left[ \gammadot \; \sign(\sigma) \right]_{n+\theta}, \\
 & \hspace{0.5cm} (\mbox{now with} \;\; \theta = 1) \nonumber \\
 & = \gammadot_{n+1} \; \sign(\sigmanp).
\end{align}
Using~(\ref{eq:timeStepApproximation}) again for $\gammadot$ gives
\begin{align}
\gammadot \approx \dst\frac{\gammanp - \gamman}{\Delta t} & 
= \gammadot_{n+\theta}, \\
 & \hspace{0.5cm} (\mbox{now with} \;\; \theta = 1) \nonumber \\
 & = \gammadotnp.
\end{align}
Thus
\be
\hstrainpdot \approx \dst\frac{\hstrainpnp - \hstrainpn}{\Delta t} 
= \dst\frac{\gammanp - \gamman}{\Delta t} \; \sign(\sigmanp).
\ee
Adopting the notation $\gammanp - \gamman = \Delta\gamma$, 
\be
\boxed{\hstrainpnp = \hstrainpn + \Delta\gamma \; \sign(\sigmanp)}.
 \label{eq:plasticStrainDiscrete}
\ee 
Analogous arguments with the hardening 
variable evolution $\alphadot = \gammadot$ lead to
\be
 \boxed{\alphanp = \alphan + \Delta\gamma}.
 \label{eq:alphaDicrete}
\ee

\subsection{Trial Elastic State and Plastic Return-Mapping}
Consider now a trial state that is purely elastic, 
with all plastic flow ``frozen'' in time. 
The incremental displacement $\Delta \vu$ creates an 
incremental strain $\Delta \hstrain \equiv \hstrainnp - \hstrainn$.
Then the trial state functions are denoted as
\begin{align}
 \sigmatrnp & = E \left( \hstrainnp - \hstrainpn \right) \equiv \sigman 
 + E \Delta \hstrain \label{eq:trialStress} \\
 \Phitrnp & = \abs{\sigmatrnp} - \left(\sigma_y 
 + H \alphan \right) \label{eq:trialYieldFunction} \\
 \hstrainptrnp & = \hstrainpn \\
 \alphatrnp & = \alphan
\end{align}
With the known initial conditions $\left\{\hstrainpn, \alphan \right\}$,
the strain is incremented $\Delta \hstrain$ 
and the resulting trial yield function is
tested for the occurrence of a strictly elastic step or an elasto-plastic step.  
Specifically,
\begin{align}
& \mbox{if} \; \Phitrnp \leq 0, \; \mbox{then} \; 
\Delta \gamma = 0 \; \mbox{and} \nonumber \\
& \left.
\begin{tabular}{rcl}
 $\sigmanp$ & $=$ & $\sigmatrnp$ \\
 $\Phinp$ & $=$ & $\Phitrnp$ \\
 $\hstrainpnp$ & $=$ & $\hstrainptrnp$ \\
 $\alphanp$ & $=$ & $\alphatrnp$ 
\end{tabular}
\right\} \; \mbox{(elastic step)}.
\end{align}
Otherwise, the return-mapping algorithm is used for the elasto-plastic step,
\begin{align}
& \mbox{if} \; \Phitrnp > 0, \; \mbox{then} \; 
\Delta \gamma > 0 \; \mbox{and} \nonumber \\
& \left.
\begin{tabular}{rcl}
 $\sigmanp$ & $=$ & $\sigmatrnp - E \Delta \gamma \; 
 \sign\left(\sigmatrnp \right)$ \\
 $\Phinp$ & $=$ & $0$ \\
 $\hstrainpnp$ & $=$ & $\hstrainpn + \Delta \gamma \; 
 \sign\left(\sigmatrnp \right)$ \\
 $\alphanp$ & $=$ & $\alphan + \Delta \gamma$ 
\end{tabular}
\right\} \; \mbox{(ep step)}.
\label{eq:elastoPlasticStep}
\end{align}
The results in~(\ref{eq:elastoPlasticStep}) are derived as follows.  First, 
find the following two useful relationships between $\sigmatrnp$ and $\sigmanp$.
\begin{align}
 \sigmanp & = E \left(\hstrainnp - \hstrainpnp \right) \\
 & = E \left(\hstrainnp - \hstrainpn + \hstrainpn - \hstrainpnp \right), \\
 & \overset{(\ref{eq:plasticStrainDiscrete}, \ref{eq:trialStress})}{=} 
 \sigmatrnp - E \Delta \gamma \; \sign(\sigmanp).
\end{align}
Thus 
\be 
\sign(\sigmanp) \left( \abs{\sigmanp} + E \Delta \gamma \right)  =  
\sign(\sigmatrnp) \abs{\sigmatrnp}.
\ee
Since $E > 0$ and $\Delta \gamma > 0$, we require
\be
\boxed{\sign(\sigmanp) = \sign(\sigmatrnp)},
\ee
and
\be
\boxed{\abs{\sigmanp} + E \Delta \gamma = \abs{\sigmatrnp}}.
\label{eq:sigmaToSigmaTrial}
\ee
Next, determine the
incremental plastic multiplier $\Delta \gamma$ using the 
discrete versions of the yield function
\begin{align}
 \Phinp 
 & \overset{(\ref{eq:yieldFunctionHardening})}{=} \abs{\sigmanp} - \left( 
 \sigma_y + H \alphanp \right) \overset{\mbox{\footnotesize set}}{=} 0, \\
 & \overset{(\ref{eq:sigmaToSigmaTrial})}{=} \abs{\sigmatrnp} - 
 E \Delta \gamma - \left(\sigma_y + H \alphan \right) \nonumber \\
 & \hspace{0.5cm} - H (\alphanp - \alphan), \\
 & \overset{(\ref{eq:trialYieldFunction}, \ref{eq:alphaDicrete})}{=} 
 \Phitrnp - E \Delta \gamma - H \Delta \gamma,  
\end{align}
gives
\be
 \boxed{\Delta \gamma = \dst\frac{\Phitrnp}{E + H}}.
\ee
Then all variables Eq.~(\ref{eq:elastoPlasticStep}) are known.  
Equation~(\ref{eq:elastoPlasticStep})$_1$, rewritten as,
\be
 \sigmanp 
 = \left[1 - \dst\frac{E \Delta \gamma}{\abs{\sigmatrnp}} \right] \sigmatrnp,
\ee
shows that the final stress 
state $\sigmanp$ may be interpreted as the {\bf projection} of the trial stress
state $\sigmatrnp$ on the yield function.  
Figure~\ref{fig:plastic_load_unload}~(b) illustrates this concept.


\begin{figure}[htb]
  \begin{center}
  %\begin{overpic}[width=\textwidth, grid, tics=5]{fig/plastic_load_unload}
  \begin{overpic}[width=\textwidth]{fig/plastic_load_unload}
  \put( 5, 75){(a)}
  \put(45, 75){(b)}
  \put(17.8,  5){$\hstrainpn$}
  \put(27,  5){$\hstrainen$}
  \put(55.5, 16){$\hstrainpn$}
  \put(62.6, 16){$\hstrainen$}
  \put(26.7, 46.3){1}
  \put(18.1, 39.3){2}
  \put(56.2, 39.3){2}
  \put(71.8, 58.6){3e}
  \put(91.3, 44.7){3}
  \put(18, 25){$E$}
  \put(56, 25){$E$}
  \put(40, 45){$\Eep$}
  \put(90, 51){$\Eep$}
  \put(16.8, 16){$\hstrainpnp$}
  \put(59,  5){$\hstrainpnp$}
  \put(23.8, 16){$\hstrainenp$}
  \put(74,  5){$\hstrainenp$}
  \put(29.5, 10){$\Delta \varepsilon_n$}
  \put(74, 10){$\Delta \varepsilon_n$}
  \put(12, 18){$O$}
  \put(50, 18){$O$}
  \put(14, 73){$\sigma$}
  \put(52, 73){$\sigma$}
  \put(44, 20){$\hstrain$}
  \put(96.5, 20){$\hstrain$}
  \put( 0, 31){$\sigmatrnp = \sigmanp$}
  \put(11, 43){$\sigma_n$}
  \put(49, 31){$\sigma_n$}
  \put(47, 66){$\sigmatrnp$}
  \put(47, 48){$\sigma_{n+1}$}
  \put(11, 37){$\sigma_y$}
  \put(49, 37){$\sigma_y$}
  \put(80, 68){\footnotesize elastic predictor}
  \put(85, 57){\footnotesize plastic corrector}
  \end{overpic}
  \end{center}
 \caption{(a) Unloading step from state 1 to state 2 
 and (b) loading step from state 2, to state 3e 
 (elastic prediction), and to state 3.}
 \label{fig:plastic_load_unload}
\end{figure}
 
\subsection{Linearization}  
  
The Taylor series expansion of real function $f(x)$ 
about a point $x=a$ is applied to 
the residual function $\vr(\vu)$ as follows,
\begin{align}
 f(x) & = f(a) \nonumber \\
      & \hspace{0cm} + f'(a) (x-a) \nonumber \\
      & \hspace{0cm} + \frac{f''(a)}{2!} (x-a)^2 \nonumber \\
      & \hspace{0cm} + \cdots \nonumber \\
      & \hspace{0cm} + \frac{f^{(n)}(a)}{n!} (x-a)^n 
\end{align}
Thus, with combining all higher order terms, abbreviated as ({\sc h.o.t.}),
\be
 \vr(\vu + \Delta \vu) 
 = \vr(\vu) 
 + \frac{\partial \vr(\vu)}{\partial \vu} \Delta \vu 
 + \mbox{\footnotesize H.O.T.},
\ee
and
\be
 \vr^{(i+1)} = \vr^{(i)} 
 + \frac{\partial \vr^{(i)}}{\partial \vu^{(i)}} \Delta \vu^{(i)} 
 + \mbox{\footnotesize H.O.T.}
\ee
This iterative approach is the groundwork for the Newton-Raphson 
iterative procedure.  At the 
initial iteration, $i=0$, the 
initial guess is set equal to the known quantity 
at the previously converged time step, 
$\vunp^{(0)} = \vun$.  Then, 
linearization about the current iterative step is taken, and 
updates are determined, as explicated in Box~\ref{box:NewtonExplicated}.
Figure~\ref{fig:newton_raphson_concept} 
shows a graphical representation of this procedure.

\begin{hboxfloat*}
\caption{Newton-Raphson iterative procedure, 
explicated for the first and second iterations.}
\label{box:NewtonExplicated}
\begin{tabular}{lll} 
   % col 1 and 
   % col 2 (multicolumn)
 \multicolumn{2}{l}{$\vunp^{(0)} = \vun$ 
initial guess as converged solution at $t_n$} 
 & % col 3
 \\
   % col 1
 & % col 2
 & % col 3
 \\
   % col 1
 $\left[ \dst\frac{-\partial \vr 
 \left(\vunp^{(0)} \right)}{\partial \vunp^{(0)}} \right]^{-1}$
 & % col 2
 $\vr \left( \vunp^{(0)} \right) = \Delta \vunp^{(0)}$
 & % col 3
 \\
   % col 1
 & % col 2
 $\implies \Delta \vunp^{(0)} + \vunp^{(0)} = \vunp^{(1)}$ 
 & % col 3
 $\implies \dst\frac{\left| \vr 
 \left( \vunp^{(1)} \right) \right|}{\abs{\vfext_{n+1}}}  
 \overset{?}{\leq} \mbox{\tt TOL}$ 
 \\
%
   % col 1
 & % col 2
 & % col 3
 \\
   % col 1
 & % col 2
 & % col 3
 \\
   % col 1
 $\left[ \dst\frac{-\partial \vr 
 \left(\vunp^{(1)} \right)}{\partial \vunp^{(1)}} \right]^{-1}$
 & % col 2
 $\vr \left( \vunp^{(1)} \right) = \Delta \vunp^{(1)}$
 & % col 3
 \\
   % col 1
 & % col 2
 $\implies \Delta \vunp^{(1)} + \vunp^{(1)} = \vunp^{(2)}$ 
 & % col 3
 $\implies \dst\frac{\left| \vr 
 \left( \vunp^{(2)} \right) \right|}{\abs{\vfext_{n+1}}}  
 \overset{?}{\leq} \mbox{\tt TOL}$ 
 \\
%
   % col 1
 & % col 2
 & % col 3
 \\
   % col 1
   $\vdots$ 
 & % col 2
   $\vdots$ 
 & % col 3 
   $\vdots$
 \\
   % col 1
 & % col 2
 & % col 3
 \\
   % col 1
   $\cdots$ 
 & % col 2
   $\cdots$ 
 & % col 3
   $\implies \dst\frac{\left| \vr 
   \left( \vunp^{(*)} \right) \right|}{\abs{\vfext_{n+1}}} 
   \leq \mbox{\tt TOL}$ 
 \\
   % col 1
 & % col 2
 & % col 3
 \\
\multicolumn{2}{l}{$\vunp = \vunp^{(*)}$ update with converged solution} & \\
\end{tabular}
\end{hboxfloat*}

\begin{figure}[htb]
  \begin{center}
  %\begin{overpic}[width=\textwidth, grid, tics=5]{fig/newton_raphson_concept}
  \begin{overpic}[width=\textwidth]{fig/newton_raphson_concept}
  \put( 1, 27){$\vfext_{n}$}
  \put( 1, 65){$\vfext_{n+1}$}
  \put( 1, 76){$\vfint$}
  \put(96, 19){$\vu$}
  \put(18.4, 23){$n$}
  \put(37.5, 73.5){$n+1$}
  \put(11, 35){$\vK^{(0)}$}
  \put(27, 53){$\vK^{(1)}$}
  \put(42, 66){$\vK^{(2)}$}
  \put(11, 10){\begin{tabular}{l} $\vu^{(0)}_{n+1}$  
  \\ $(\uparrow \; \equiv \; \vu_n)$ \end{tabular}}
  \put(24, 11.5){$\vu^{(1)}_{n+1}$}  
  \put(39, 11.5){$\vu^{(2)}_{n+1} \cdots$}    
  \put(47, 10){\begin{tabular}{l} $\vu^{(*)}_{n+1}$  
  \\ $(\downarrow \; \equiv \; \vunp)$ \end{tabular}}
  \put( 5, 17){$O$}
  \put(16.5,  3){$\Delta\vunp^{(0)}$}    
  \put(30,  3){$\Delta\vunp^{(1)}$}      
  \put(82, 35){$\vr^{(0)}$}
  \put(74, 53){$\vr^{(1)}$}
  \put(60.5, 61.5){$\vdots \;\;\;\;\; \;\;  \vr^{(2)}$}
  \put(57, 66){$\vr^{(*)} \approx \ve{0}$}
  \end{overpic}
  \end{center}
  \caption{Newton-Raphson iterative procedure, shown graphically:  The
  first iteration $K^{(0)}$ approximates the curve as tangent line at $n$.
  Successive tangents are constructed $K^{(i)}$ until the root at
  $n+1$ is found.}
  \label{fig:newton_raphson_concept}
\end{figure}




\begin{hboxfloat*}
 \caption{Predictor/corrector time stepping algorithm 
 applied to elasto-plastic equations of motion.}
 \label{box:defimp}
%\renewcommand{\baselinestretch}{2.0}
%\small % make font smaller
\normalsize 
\begin{center}
{\bf Time step loop} ($n=0$, {\tt WHILE} $n \leq $ {\tt NSTEPS})
\begin{enumerate}
 	\item $i=0$, {\tt CONVERGED = FALSE} 
	\item $\vunpi = \vun$ (prediction)
	%
	\vspace{0.5cm}
      \item {\bf Iteration loop} ({\tt WHILE} $i \leq \;${\tt IMAX 
      AND CONVERGED IS FALSE})
      \begin{itemize}
      	%
		\vspace{0.25cm}
		\item {\em if} \hspace{0.25cm} 
			$\dst\frac{\left| \vr 
			\left( \vunpi \right) \right|}{\abs{\vfext_{n+1}}}  
			\leq$ {\tt TOL}
			\hspace{0.25cm} $\abs{\vfext_{n+1}} \neq 0$
		\vspace{0.125cm}
		\item {\em then}
		\begin{enumerate}
			\item {\tt CONVERGED = TRUE}
			\item $\vunp^{(*)} = \vunpi$
			\item $\hstrainnp^{\mbox{\footnotesize p}\;(*)} 
			      = \hstrainpnpi$
			\item $\alphanp^{(*)} = \alphanpi$
		\end{enumerate}
		\item {\em else}
		\vspace{0.25cm}
		\begin{enumerate}
			\item $\left[ 
			\dst\frac{-\partial 
			\vr \left(\vunpi \right)}{\partial \vunpi} 
			\right]^{-1} \vr \left( \vunpi \right) 
			= \Delta \vunpi$ (form tangent and solve)
      	%
		\vspace{0.25cm}
			\item $\vunpip = \vunpi 
			+ \Delta \vunpi$ (iteration update)
			\item $i \leftarrow i+1$
		\end{enumerate}
	\end{itemize}
	%
	\vspace{0.5cm}
      \item $\vunp = \vunp^{(*)}$ (time step update)
      \item $\hstrainpnp = \hstrainnp^{\mbox{\footnotesize p}\;(*)}$
      \item $\alphanp = \alphanp^{(*)}$
      \item $n \leftarrow n+1$
      \end{enumerate}  
\end{center}
\end{hboxfloat*}











  














